\section{Proofs for the table of outcomes}\label{app:examples}

This appendix proves the rows $P_d$ and $\Theta_\varepsilon$ of
Table~\ref{tab:outcomes}. In both rows $Q$ is symmetric and the start has full
support. So the admissible faces are the connected faces
(Lemma~\ref{lem:admissible}), and the winners at $\tau$ are the connected faces
that minimize $\Phi_\tau(F)=\tau\lambda_F+\lvert F\rvert-1$
(Theorem~\ref{thm:hull}). We use three facts about a connected face $F$. First,
$-Q_F$ is symmetric, $\lambda_F$ is its least eigenvalue, and an eigenvector of
$-Q_F$ with positive entries belongs to $\lambda_F$ (the second fact at the
start of Appendix~\ref{app:torus}). Second, $-Q_F$ is the Laplacian of the
graph on $F$, which is positive semidefinite, plus the diagonal matrix of the
rates out of $F$. So $\lambda_F$ is the least value of
\begin{equation}\label{eq:rayleigh}
 \sum_{\{i,j\}\subseteq F}q_{ij}(f_i-f_j)^2+\sum_{i\in F}q_i^{\rm out}(F)f_i^2
\end{equation}
over unit vectors $f\in\R^F$. Third, if $F'\subsetneq F$ is nonempty, then
$\lambda_{F'}>\lambda_F$ (Proposition~\ref{prop:rate}(d)).

Each proof has two steps, namely the exit rates and then the lower-left boundary
of the hull. For the second step we use the following lemma. It says that a
chain of points with increasing switch values is its own lower-left boundary.

\begin{lemma}[a chain of lines]\label{lem:chain}
Let $x_0>x_1>\dots>x_m\ge0$ and $y_0<y_1<\dots<y_m$, and put
$\ell_i(\tau)=\tau x_i+y_i$ and $t_i=(y_i-y_{i-1})/(x_{i-1}-x_i)$ for
$1\le i\le m$. Put $t_0=0$ and $t_{m+1}=\infty$, and suppose that
$t_1<t_2<\dots<t_m$. Then for $1\le j\le m$ and every $\tau$
\begin{equation}\label{eq:chain}
 \ell_j(\tau)-\ell_{j-1}(\tau)=(x_{j-1}-x_j)(t_j-\tau).
\end{equation}
Let $0\le i\le m$ and let $\tau$ be finite with $t_i\le\tau\le t_{i+1}$. Then
$\ell_0(\tau)\ge\dots\ge\ell_i(\tau)\le\dots\le\ell_m(\tau)$. If also
$t_i<\tau<t_{i+1}$, then all these inequalities are strict, so
$\ell_i(\tau)<\ell_j(\tau)$ for every $j\ne i$.
\end{lemma}

\begin{proof}
The identity~\eqref{eq:chain} is the definition of $t_j$. If $j\le i$, then
$t_j\le t_i\le\tau$, so $\ell_j(\tau)\le\ell_{j-1}(\tau)$. If $j>i$, then
$\tau\le t_{i+1}\le t_j$, so $\ell_j(\tau)\ge\ell_{j-1}(\tau)$. If
$t_i<\tau<t_{i+1}$, these inequalities are strict.
\end{proof}

\subsection{The path}\label{app:paths}

Let $P_d$ be the path with states $1,\dots,d$ and $q_{i,i+1}=q_{i+1,i}=1$ for
$1\le i<d$. An \emph{end-arc} of $k<d$ states is one of the two faces
$\{1,\dots,k\}$ and $\{d-k+1,\dots,d\}$. For $k\ge1$ put
\[
 \lambda'_k=2-2\cos\frac{\pi}{2k+1},\qquad
 \tau'_k=\frac1{\lambda'_{k-1}-\lambda'_k}\ (k\ge2),\qquad
 D'_k=k+\lambda'_k\tau'_k\ (k\ge2).
\]

\begin{lemma}[the numbers $D'_k$]\label{lem:Dprime}
\begin{enumerate}[(a)]
\item The sequence $\lambda'_k$ is positive, strictly decreasing and strictly
convex in $k\ge1$. Hence $0<\tau'_2<\tau'_3<\cdots$. Also
$\tau'_2=(1+\sqrt5)/2$.
\item $\frac32k-\frac12<D'_k<\frac32k$ for every $k\ge2$.
\end{enumerate}
\end{lemma}

\begin{proof}
(a) Let $g(x)=2-2\cos(\pi/x)$, so that $\lambda'_k=g(2k+1)$. For $x\ge2$,
\[
 g'(x)=-\frac{2\pi\sin(\pi/x)}{x^2}<0,\qquad
 g''(x)=2\pi\Bigl(\frac{\pi\cos(\pi/x)}{x^4}+\frac{2\sin(\pi/x)}{x^3}\Bigr)>0 .
\]
So $\lambda'_{k-1}-\lambda'_k>\lambda'_k-\lambda'_{k+1}>0$ for $k\ge2$, and
the $\tau'_k$ increase. Since $\lambda'_1=1$ and
$\lambda'_2=2-2\cos\frac\pi5=(3-\sqrt5)/2$, we get
$\tau'_2=2/(\sqrt5-1)=(1+\sqrt5)/2$.

(b) Put $a=\pi/(2(2k-1))$ and $b=\pi/(2(2k+1))$, so that $0<b<a\le\pi/6$. Since
$2-2\cos2x=4\sin^2x$, the ratio $\rho=\lambda'_k/\lambda'_{k-1}$ equals
$(\sin b/\sin a)^2$. Also $D'_k-k=\rho/(1-\rho)$, which increases with $\rho$.
So the claim is
\[
 \frac{k-1}{k+1}<\rho<\frac k{k+2}.
\]
For the lower bound, $\sin x/x$ decreases on $(0,\pi]$, so
$\sin b/\sin a>b/a=(2k-1)/(2k+1)$. Then $\rho>(2k-1)^2/(2k+1)^2>(k-1)/(k+1)$,
because $(2k-1)^2(k+1)-(k-1)(2k+1)^2=2$.

For the upper bound write $s(x)=\sin x/x$ and
$y=(8k-2)/((k+2)(2k-1)^2)$. Then
\[
 \rho=\Bigl(\frac{2k-1}{2k+1}\Bigr)^2\Bigl(\frac{s(b)}{s(a)}\Bigr)^2
 \qquad\text{and}\qquad
 \frac k{k+2}=\Bigl(\frac{2k-1}{2k+1}\Bigr)^2(1+y),
\]
the second because $k(2k+1)^2-(k+2)(2k-1)^2=8k-2$. So we must show that
$s(b)/s(a)<\sqrt{1+y}$. Now $0<y<1$, because
$(k+2)(2k-1)^2\ge12(2k-1)>8k-2$, and $\sqrt{1+y}\ge1+\frac y3$ for
$0\le y\le3$. So it suffices that $s(b)-s(a)<\frac y3s(a)$. We use
$1-\frac{x^2}6\le s(x)\le1-\frac{x^2}6+\frac{x^4}{120}$ for $x>0$. This gives
$s(a)\ge1-\frac{\pi^2}{216}>0.954$ and
\[
 s(b)-s(a)\le\frac{a^2-b^2}6+\frac{b^4}{120}
 =\frac{\pi^2k}{3(4k^2-1)^2}+\frac{\pi^4}{1920(2k+1)^4}.
\]
So it suffices that this is less than $0.318\,y$. Multiply by $(2k-1)^2$ and
use $k/(2k+1)^2<1/(4k)$ and $(2k-1)^2/(2k+1)^4<1/(4k^2)$. It then suffices
that
\[
 \frac{\pi^2}{12k}+\frac{\pi^4}{7680k^2}<0.318\,\frac{8k-2}{k+2}.
\]
The left side decreases in $k$ and the right side, which is
$0.318(8-\frac{18}{k+2})$, increases. At $k=2$ they are less than $0.415$ and
more than $1.11$.
\end{proof}

\begin{proposition}[the path]\label{prop:paths}
Let $d\ge3$, let $Q$ be the path $P_d$ with unit rates, and let the start have
full support. Put $K=\lfloor2d/3\rfloor$ and
$\tau_{\rm full}=\lceil d/3\rceil/\lambda'_K$.
\begin{enumerate}[(a)]
\item The connected faces are the sets of consecutive states. An end-arc of
$k<d$ states has exit rate $\lambda'_k$. A set of $k$ consecutive states that
contains neither $1$ nor $d$ has exit rate $2-2\cos\frac\pi{k+1}>\lambda'_k$.
\item We have $2\le K\le d-1$ and
$\tau'_2<\dots<\tau'_K<\tau_{\rm full}$.
\item The winners are the states $1$ and $d$ for $0<\tau<\tau'_2$, the two
end-arcs of $k$ states for $\tau'_k<\tau<\tau'_{k+1}$ and $2\le k\le K-1$, the
two end-arcs of $K$ states for $\tau'_K<\tau<\tau_{\rm full}$, and $P_d$ alone
for $\tau>\tau_{\rm full}$. No other face wins for any $\tau>0$.
\end{enumerate}
\end{proposition}

In words, the mode grows from an end of the path, one state at a time, until
it covers two thirds of the path, and then it jumps to the whole path. The
first step is at the golden ratio.

\emph{Idea of the proof.} Growing an end-arc by one state lowers its exit rate
by less and less, since $\lambda'_k$ is convex. So the end-arcs alone win in
the order of their size, and the arc of $k$ states starts to win at
$\tau'_k$. The whole path enters when its exponent $d-1$ falls below the least
exponent of the arcs. At $\tau'_k$ that least exponent is $D'_k-1$, so the
arcs of $k$ states win before the whole path exactly when $D'_k<d$.
Lemma~\ref{lem:Dprime}(b) puts $D'_k$ just below $\frac32k$, and this gives
$k\le2d/3$.

\begin{proof}
(a) A connected face of a path is a set of consecutive states. Let $F$ have
$k<d$ states. If $F=\{1,\dots,k\}$, then $-Q_F$ is tridiagonal with
off-diagonal entries $-1$ and diagonal $(1,2,\dots,2)$. Put
$\vartheta=\pi/(2k+1)$ and $r_j=\cos((2j-1)\vartheta/2)$ for $1\le j\le k$.
Then $r_j>0$, since $(2j-1)\vartheta/2<\pi/2$. With $r_0=r_1$ and
$r_{k+1}=\cos(\pi/2)=0$ we have $r_{j-1}+r_{j+1}=2\cos\vartheta\,r_j$ for
$1\le j\le k$, and this says that $-Q_Fr=(2-2\cos\vartheta)r$. So
$\lambda_F=\lambda'_k$, and the same holds for $\{d-k+1,\dots,d\}$ by the
reflection $i\mapsto d+1-i$. If $F$ contains neither $1$ nor $d$, then
$-Q_F=2I-A$ with $A$ the adjacency matrix of a path on $k$ states. Now
$r_j=\sin(j\pi/(k+1))$ is a positive eigenvector with eigenvalue
$2-2\cos\frac\pi{k+1}$, and this is larger than $\lambda'_k$ because
$\frac\pi{k+1}>\frac\pi{2k+1}$.

(b) Since $d\ge3$ we have $2\le K\le d-1$. The $\tau'_k$ increase by
Lemma~\ref{lem:Dprime}(a). Note that $d-K=\lceil d/3\rceil$, so
$\tau_{\rm full}=(d-K)/\lambda'_K$. By Lemma~\ref{lem:Dprime}(b),
$D'_K<\frac32K\le d$, that is $\lambda'_K\tau'_K<d-K$. So
$\tau'_K<\tau_{\rm full}$.

(c) Write $\phi_k(\tau)=\tau\lambda'_k+k-1$ for the exponent of an end-arc of
$k$ states, $1\le k\le d-1$. First we apply Lemma~\ref{lem:chain} to the
points $(\lambda'_k,k-1)$ for $1\le k\le K$ followed by $(0,d-1)$, the point
of $P_d$. Its switch values are $\tau'_2,\dots,\tau'_K$ and
$(d-K)/\lambda'_K=\tau_{\rm full}$, which increase by (b). So on each of the
intervals in (c) the face named there has a smaller exponent than the other
end-arcs of at most $K$ states and than $P_d$, and by (a) also than every set
of at most $K$ consecutive states that is not an end-arc.

It remains to show that a face $F$ of $k$ consecutive states with
$K<k<d$ never wins. By (a) its exponent is at least $\phi_k(\tau)$. We show that
$\phi_k(\tau)>\min(\phi_K(\tau),d-1)$ for every $\tau>0$. Now we apply
Lemma~\ref{lem:chain} to the points $(\lambda'_j,j-1)$ for $1\le j\le d-1$,
with switch values $\tau'_2<\dots<\tau'_{d-1}$. If $\tau<\tau'_{K+1}$, then
$\tau<\tau'_j$ for $K<j\le k$, and~\eqref{eq:chain} gives
$\phi_k(\tau)>\phi_{k-1}(\tau)>\dots>\phi_K(\tau)$. If $\tau\ge\tau'_{K+1}$, take $j$
with $K+1\le j\le d-1$ and $\tau'_j\le\tau$, and with $\tau\le\tau'_{j+1}$ if
$j<d-1$. Then
\[
 \phi_k(\tau)\ge \phi_j(\tau)\ge \phi_j(\tau'_j)=D'_j-1>\tfrac32j-\tfrac32
 \ge\tfrac32K\ge d-1 .
\]
The first step is Lemma~\ref{lem:chain}, the second holds since $\phi_j$
increases, the third is the definition of $D'_j$, the fourth is
Lemma~\ref{lem:Dprime}(b), the fifth is $j\ge K+1$, and the last holds because
$3K\ge2d-2$. So $F$ never wins. The winners at $\tau$ are the admissible faces with the
least exponent, and this gives (c).
\end{proof}

At a switch value the faces on both sides tie. On $P_4$ we have $K=2$ and
$D'_3=4.077>4$, so the end-arc of three states never wins. Its point lies
only $\lambda'_3-\frac12\lambda'_2=0.0071$ to the right of the edge from the
end-arcs of two states to $P_4$. The distance is small because the lower bound
in Lemma~\ref{lem:Dprime}(b) is nearly attained.

\subsection{Two triangles}\label{app:triangles}

Let $\varepsilon>0$. The graph $\Theta_\varepsilon$ has states $1,\dots,6$,
unit rates on the triangles $A=\{1,2,3\}$ and $B=\{4,5,6\}$, and one more link
between $3$ and $4$ of rate $\varepsilon$. We call $1,2,5,6$ the inner states
and $\{1,2\}$, $\{5,6\}$ the inner edges. The graph has the symmetries
$1\leftrightarrow2$, $5\leftrightarrow6$ and the mirror map $i\mapsto7-i$. Put
\begin{align*}
 p_A(x)&=x^2-(3+\varepsilon)x+\varepsilon,\\
 p_4(x)&=x^3-(5+2\varepsilon)x^2+(6+6\varepsilon)x-2\varepsilon,\\
 p_5(x)&=x^3-(4+2\varepsilon)x^2+(3+4\varepsilon)x-\varepsilon,
\end{align*}
and $\varepsilon_0=(64-\sqrt{46})/60=0.95363$,
$\varepsilon_1=(64+\sqrt{46})/60=1.17971$, the roots of
$120\varepsilon^2-256\varepsilon+135$.

\begin{proposition}[two triangles joined by a link]\label{prop:triangles}
Let the start have full support.
\begin{enumerate}[(a)]
\item $\lambda_A=\lambda_B=\frac12\bigl(3+\varepsilon-\sqrt{\varepsilon^2+2\varepsilon+9}\bigr)$
is the smaller root of $p_A$, and $0<\lambda_A<1$.
\item Among the connected faces of one, two and three states the least exit
rates are $2$, $1$ and $\lambda_A$. They are attained exactly by the inner
states, by the inner edges and by $A$ and $B$. Put
$\lambda_4=\lambda_{\{1,2,3,4\}}$ and $\lambda_5=\lambda_{\{2,\dots,6\}}$. Every
connected face of four states has exit rate at least $\lambda_4$, and every
connected face of five states has exit rate $\lambda_5$. Here $\lambda_4$ and
$\lambda_5$ are the least roots of $p_4$ and $p_5$, and
$\lambda_5<\lambda_4<\lambda_A$.
\item $\lambda_5>\frac12\lambda_4$ for every $\varepsilon>0$, and
$\lambda_4>\frac23\lambda_A$ if and only if $\varepsilon<\varepsilon_1$.
\item Let $\varepsilon<\varepsilon_1$, and put
\[
 \tau_A=\frac1{1-\lambda_A}=\frac{1+\varepsilon+\sqrt{\varepsilon^2+2\varepsilon+9}}4,
 \qquad
 \tau^*=\frac3{\lambda_A}=\frac{3\bigl(3+\varepsilon+\sqrt{\varepsilon^2+2\varepsilon+9}\bigr)}{2\varepsilon}.
\]
Then $1<\tau_A<\tau^*$. The winners are the inner states for $0<\tau<1$, the
inner edges for $1<\tau<\tau_A$, the triangles $A$ and $B$ for
$\tau_A<\tau<\tau^*$, and $[6]$ alone for $\tau>\tau^*$. No other face wins
for any $\tau>0$. As $\varepsilon\to0$,
$\tau_A=1+\frac\varepsilon3+O(\varepsilon^2)$ and
$\tau^*=\frac9\varepsilon+2+O(\varepsilon)$.
\item Let $\varepsilon>\varepsilon_1$. Then at $\tau=3/\lambda_A$, where the
triangles and $[6]$ have the same exponent, the face $\{1,2,3,4\}$ has a
smaller exponent. So the list in (d) fails.
\end{enumerate}
\end{proposition}

In words, for a weak link each triangle is a community that the chain rarely
leaves. The mode sits on a triangle from $\tau\approx1$ until
$\tau\approx9/\varepsilon$, and then it spreads over all six states at once.
Once the link is stronger than $\varepsilon_1$ this list fails, because a
triangle together with the far end of the link beats both the triangles and
$[6]$ at the value of $\tau$ where they tie.

\emph{Idea of the proof.} The exit rates of small faces are explicit, and
those of the faces with four and five states are roots of cubics. The points
of the inner states, the inner edges, the triangles and $[6]$ form a chain as
in Lemma~\ref{lem:chain}. The points of the faces with four and five states
must lie strictly to the right of the edge from a triangle to $[6]$, and this
is what (c) says. We prove (c) by continuity in
$\varepsilon$. A change of sign needs a common root of two polynomials, hence a
zero of a resultant, and the resultants are explicit polynomials in
$\varepsilon$.

\begin{proof}
(a) On the vectors $(a,a,b)$ the matrix $-Q_A$ acts as
$(a,b)\mapsto(a-b,\,-2a+(2+\varepsilon)b)$, with characteristic polynomial
$p_A$, and $(1,-1,0)$ is an eigenvector with eigenvalue $3$. As
$p_A(0)=\varepsilon>0$ and $p_A(1)=-2<0$, the smaller root of $p_A$ lies in
$(0,1)$ and the other root is larger than $1$. So the smaller root is
$\lambda_A$. The mirror map gives $\lambda_B=\lambda_A$.

(b) By~\eqref{eq:rayleigh}, $\lambda_F\ge1$ whenever every state of $F$ has
rate at least $1$ out of $F$. We go through the connected faces by size. The
states $3$ and $4$ have exit rate $2+\varepsilon$, and the inner states have
$2$. The connected pairs are the inner edges, with $\lambda=1$, the pair
$\{3,4\}$, with $\lambda=2$, and $\{1,3\}$, $\{2,3\}$ and their mirror images,
with $\lambda=\frac12(4+\varepsilon-\sqrt{\varepsilon^2+4})>1$. The connected
faces of three states are $A$, $B$ and the sets $\{a,3,4\}$ and $\{3,4,b\}$
with $a\in\{1,2\}$ and $b\in\{5,6\}$, because a connected face that meets both
triangles contains $3$ and $4$. In each of these four sets every state has
rate at least $1$ out of the set, so $\lambda\ge1>\lambda_A$. A connected face
of four states meets both triangles, so it contains $3$ and $4$. So it is
$\{1,2,3,4\}$, its mirror image $\{3,4,5,6\}$, or a set $\{a,3,4,b\}$, in which every state has rate $1$ out
of the set, so that $\lambda\ge1$. A connected face of five states is the
complement of an inner state (the complements of $3$ and of $4$ are
disconnected), and these four faces are images of $\{2,\dots,6\}$
under the symmetries. By the third fact above,
$\lambda_5<\lambda_{\{3,4,5,6\}}=\lambda_4<\lambda_A<1$, since
$\{2,\dots,6\}\supsetneq\{3,4,5,6\}$ and $\{1,2,3,4\}\supsetneq A$.

For the polynomials, note that $(1,-1,0,0)$ is an eigenvector of $-Q_F$ for
$F=\{1,2,3,4\}$ with eigenvalue $3$. On the vectors $(a,a,b,c)$ the matrix
acts as
$(a,b,c)\mapsto(a-b,\,-2a+(2+\varepsilon)b-\varepsilon c,\,-\varepsilon b+(2+\varepsilon)c)$,
with characteristic polynomial
$(x-1)(x-2)(x-2-2\varepsilon)-2(x-2-\varepsilon)=p_4(x)$. So the eigenvalues
of $-Q_F$ are $3$ and the roots of $p_4$. As $\lambda_4<1$, it is the least
root of $p_4$. In the same way, for $F=\{2,\dots,6\}$ the vector that is $1$
at $5$, $-1$ at $6$ and $0$ elsewhere has eigenvalue $3$, and on the vectors
that agree at $5$ and $6$ the characteristic polynomial is $(x-3)p_5(x)$. So
$\lambda_5$ is the least root of $p_5$. (Both characteristic polynomials are
also checked by computer algebra, see Section~\ref{app:examples-checks}.)

(c) Put $D_1=\lambda_5-\frac12\lambda_4$ and
$D_2=\lambda_4-\frac23\lambda_A$. Both are continuous in $\varepsilon>0$,
because the least eigenvalue of a symmetric matrix is continuous in its
entries. If $D_1=0$, then $x=\lambda_5$ is a common root of $p_5(x)$ and
$p_4(2x)$. If $D_2=0$, then $x=\lambda_4$ is a common root of $p_4(x)$ and
$p_A(\frac32x)$. A common root makes the resultant vanish, and an exact
computation gives
\begin{align*}
 \operatorname{Res}_x\bigl(p_5(x),p_4(2x)\bigr)
 &=-8\varepsilon^3(16\varepsilon^2+12\varepsilon+9),\\
 \operatorname{Res}_x\bigl(p_A(\tfrac32x),p_4(x)\bigr)
 &=-\tfrac1{16}\varepsilon^2(120\varepsilon^2-256\varepsilon+135).
\end{align*}
So $D_1$ has no zero in $\varepsilon>0$, and $D_2$ can vanish only at
$\varepsilon_0$ and at $\varepsilon_1$.

We find the signs for small $\varepsilon$. By (a) and (b),
$0<\lambda_5<\lambda_4<\lambda_A<\varepsilon$. Here $\lambda_5>0$ because
$\lambda_5\ge0$ by~\eqref{eq:rayleigh} and $p_5(0)=-\varepsilon\ne0$, and
$\lambda_A<\varepsilon$ because $p_A(\varepsilon)=-2\varepsilon<0$. At $\varepsilon=0$ the cubics are
$x(x-2)(x-3)$ and $x(x-1)(x-3)$, with the simple root $0$. So the implicit
function theorem gives one root $\frac\varepsilon3+O(\varepsilon^2)$ of each
cubic near $0$, and no other root near $0$. So
$\lambda_4=\frac\varepsilon3+O(\varepsilon^2)$ and
$\lambda_5=\frac\varepsilon3+O(\varepsilon^2)$, and (a) gives
$\lambda_A=\frac\varepsilon3+O(\varepsilon^2)$. Hence
$D_1=\frac\varepsilon6+O(\varepsilon^2)$ and
$D_2=\frac\varepsilon9+O(\varepsilon^2)$ are positive for small
$\varepsilon$. So $D_1>0$ for every $\varepsilon>0$, and $D_2>0$ on
$(0,\varepsilon_0)$.

Next, $D_2<0$ for large $\varepsilon$. Indeed $\lambda_A\to1$ as
$\varepsilon\to\infty$ by (a). On the other hand take $f=(a,a,c,c)$
in~\eqref{eq:rayleigh} for $\{1,2,3,4\}$. The value is $2(a-c)^2+2c^2$ and
$\lvert f\rvert^2=2a^2+2c^2$, so
$\lambda_4\le\min_{a,c}((a-c)^2+c^2)/(a^2+c^2)=(3-\sqrt5)/2<\frac12$ for every
$\varepsilon$.

Now we show that $D_2>0$ at $\varepsilon_0$. Put $x_0=\frac7{40}$. Then
\[
 p_4(x_0)=\frac{57743}{64000}-\frac{809}{800}\varepsilon,\qquad
 p_A(\tfrac32x_0)=\frac{59}{80}\varepsilon-\frac{4599}{6400}
\]
are both negative for $0.8923<\varepsilon<0.9743$, and this interval contains
$\varepsilon_0$. For $0\le x\le\frac12$ we have
$p_4'(x)=3x^2-2(5+2\varepsilon)x+6+6\varepsilon\ge1+4\varepsilon>0$. So $p_4$
increases on $[0,x_0]$ and is negative there, and its positive root
$\lambda_4$ is larger than $x_0$. Also $p_A(\frac32x_0)<0$ puts $\frac32x_0$
between the roots of $p_A$, so $\lambda_A<\frac32x_0$. Hence
$\lambda_4>x_0>\frac23\lambda_A$ at $\varepsilon_0$.

So $D_2>0$ on $(0,\varepsilon_0]$. It has no zero in
$(\varepsilon_0,\varepsilon_1)$, so $D_2>0$ on $(0,\varepsilon_1)$. It has no
zero in $(\varepsilon_1,\infty)$ and is negative for large $\varepsilon$, so
$D_2<0$ on $(\varepsilon_1,\infty)$. By continuity $D_2=0$ at
$\varepsilon_1$. This proves (c).

(d) As $p_A(\frac34)=\frac\varepsilon4-\frac{27}{16}<0$ for
$\varepsilon<\frac{27}4$, we have $\lambda_A<\frac34$. So $1-\lambda_A>0$ and
$\lambda_A<3(1-\lambda_A)$, that is $1<\tau_A<\tau^*$. Apply
Lemma~\ref{lem:chain} to the points $(2,0)$, $(1,1)$, $(\lambda_A,2)$ and
$(0,5)$ of the inner states, the inner edges, the triangles and $[6]$. The
switch values are $1$, $\tau_A$ and $\tau^*$. So on each interval in (d) the
faces named there have a smaller exponent than the other three kinds, and by (b)
than every other connected face of at most three states. Let $F$ be a
connected face of four or five states. By (b) and (c), for every $\tau>0$
\begin{align*}
 \Phi_\tau(F)&\ge\tau\lambda_4+3>\tfrac23(\tau\lambda_A+2)+\tfrac13\cdot5
 &&\text{if }\lvert F\rvert=4,\\
 \Phi_\tau(F)&=\tau\lambda_5+4>\tfrac13(\tau\lambda_A+2)+\tfrac23\cdot5
 &&\text{if }\lvert F\rvert=5,
\end{align*}
where the second line uses $\lambda_5>\frac12\lambda_4>\frac13\lambda_A$. Each
right side is at least the smaller of the exponents $\tau\lambda_A+2$ of a
triangle and $5$ of $[6]$. So $F$ never wins, and this gives the winners. The closed forms follow from (a) by rationalizing, since
$(\varepsilon^2+2\varepsilon+9)-(1+\varepsilon)^2=8$ and
$(3+\varepsilon)^2-(\varepsilon^2+2\varepsilon+9)=4\varepsilon$. The
expansions follow from
$\sqrt{\varepsilon^2+2\varepsilon+9}=3+\frac\varepsilon3+O(\varepsilon^2)$.

(e) At $\tau=3/\lambda_A$ the triangles and $[6]$ have the exponent $5$, and
$\{1,2,3,4\}$ has $3\lambda_4/\lambda_A+3<5$ by (c).
\end{proof}

At $\varepsilon=1$, for example, $\tau_A=(1+\sqrt3)/2$ and
$\tau^*=6+3\sqrt3$.

\subsection{What was checked by computation}\label{app:examples-checks}

The proofs above use computer algebra in one place only, namely the two
resultants in the proof of Proposition~\ref{prop:triangles}(c). The script
\texttt{verify\_examples.py} computes them in sympy. It also expands the
characteristic polynomials of $-Q_F$ for $F=A$, $\{1,2,3,4\}$ and
$\{2,\dots,6\}$ and finds $(x-3)p_A$, $(x-3)p_4$ and $(x-3)^2p_5$, and it
recomputes $p_4(\frac7{40})$ and $p_A(\frac{21}{80})$ in rational arithmetic.
The rest of its checks for these two rows are in floating point, and no proof
rests on them. For $3\le d\le9$ it computes the exit rates of all connected
faces of $P_d$, compares those of the end-arcs with $\lambda'_k$, and finds
the sizes of the winners as in Proposition~\ref{prop:paths}(c). With 50
digits it confirms $\frac32k-\frac12<D'_k<\frac32k$ for $2\le k\le2000$. For
$\varepsilon=0.3$, $1$ and $1.1$ it computes the exit rates of all 63 faces of
$\Theta_\varepsilon$ and finds the least exit rates, the winners and the
switch values of Proposition~\ref{prop:triangles}. For $\varepsilon=1.5$, $3$,
$4.4$, $4.6$ and $8$ it finds that the faces $\{1,2,3,4\}$ and $\{3,4,5,6\}$
win on an interval of $\tau$, which agrees with
Proposition~\ref{prop:triangles}(e).
